\documentclass[
10pt, % Main document font size
a4paper, % Paper type, use 'letterpaper' for US Letter paper
oneside, % One page layout (no page indentation)
%twoside, % Two page layout (page indentation for binding and different headers)
headinclude,footinclude, % Extra spacing for the header and footer
BCOR5mm, % Binding correction
]{scrartcl}
\input{stru/structure.tex}
\hyphenation{Fortran hy-phen-ation} 

\title{\normalfont{FTML practical session 13}\\
} 
\author{\spacedlowsmallcaps{}
} 
\date{\today}
\begin{document}
\renewcommand{\sectionmark}[1]{\markright{\spacedlowsmallcaps{#1}}}
\lehead{\mbox{\llap{\small\thepage\kern1em\color{halfgray} \vline}\color{halfgray}\hspace{0.5em}\rightmark\hfil}} 
\pagestyle{scrheadings}
\maketitle 
\setcounter{tocdepth}{1} 

\begin{figure}[htpb]
    \centering
    \includegraphics[width=1\textwidth]{Hoeffding_epsilon_0_100}
    \caption{Simulation of Hoeffding's inequality}
    \label{fig:hoeffding}
\end{figure}

\tableofcontents

\pagebreak

\section{\large\color{Blue}Hoeffding's inequality}
\label{sec:Hoeffding}

The following result will be useful in order to proove the bound on the estimation error in section \ref{sec:estimation_error}


\begin{theorem}{Hoeffding's inequality}

    Let $(X_i)_{1\leq i\leq n}$ be $n$ i.i.d real random variables such that $\forall i\in [1, n]$, $X_i\in [a,b]$ and $ E(X_i)= \mu \in \mathbb{R} $.  Let $	\bar{\mu} = \frac{1}{n} \sum^{n}_{i=1} X_i$.

    Then $\forall \epsilon >0$,

    \begin{equation}
	P\Big(| \bar{\mu}-\mu|\geq \epsilon \Big) \leq 2 \exp\Big( -\frac{2n\epsilon^2}{(b-a)^2} \Big)
    \end{equation}
\end{theorem}

{\color{Purple}{ Run a simulation that allows to visualize Hoeffding's inequality, with a random
variable of your choice, like in figures \ref{fig:hoeffding_01} and
\ref{fig:hoeffding_001}, where a Bernoulli variable of parameter $p=1/2$ is
used. }}

\begin{figure}[H]
    \centering
    \includegraphics[width=0.8\textwidth]{Hoeffding_epsilon_0_100}
    \caption{Hoeffding's inequality with a Bernoulli variable of parameter $p=1/2$ and $\epsilon=0.1$}
    \label{fig:hoeffding_01}
\end{figure}

\begin{figure}[H]
    \centering
    \includegraphics[width=0.8\textwidth]{Hoeffding_epsilon_0_010}
    \caption{Hoeffding's inequality with a Bernoulli variable of parameter $p=1/2$ and $\epsilon=0.01$}
    \label{fig:hoeffding_001}
\end{figure}

\pagebreak

\section{\large\color{Blue}Bound on the estimation error}
\label{sec:estimation_error}

We consider a usual supervised learning setting, and the  a space of functions $F$ in which we choose our estimators. The dataset contains $n$ samples, the empircial risk is noted $R_n$ and the real risk $R$. If $f_n$ is the empircial risk minimizer, and $f_a$ the optimal estimator in $F$, we have seen this result during the lectures:

\begin{equation}
    \label{eq:est_sup}
	0\leq R(f_n) - R(f_a)\leq2\sup_{h\in F}|R(h)-R_n(h)|
\end{equation}

Our objective is to bound equation \ref{eq:est_sup}. We make the following additional hypotheses:
\begin{itemize}
	\item $F$ is finite, with $|F|$ elements.
	\item The loss $l$ is uniformly bounded: $l( \hat{y}, y)\in [a, b]$ with $a$ and $b$ real numbers.
\end{itemize}

We will also use the following result:


\begin{proposition}{Boole's inequality}

    Let $A_1, A_2, \dots, $ be  acountable set of events of a probability space $\{\Omega, \mathcal{F}, P)$.

    Then, 

    \begin{equation}
	P\Big( \cup_{i\geq 1}A_i \Big) \leq \sum_{i\geq 1} P(A_i)
    \end{equation}
    
\end{proposition}

\textbf{Step 1]} Using \ref{eq:est_sup}, we have that:

    \begin{equation}
	P\Big(  R( f_n)-R(f_a)\geq t\Big) \leq P\Big(2 \sup_{h\in F}|R(h)-R_n(h)|\geq t\Big)
    \end{equation}


\textbf{Step 2]}

Show that

    
    \begin{equation}
	P\Big(2 \sup_{h\in F}|R(h)-R_n(h)|\geq t\Big) \leq \sum_{h\in F} P\big( 2 |R(h)-R_n(h)|\geq t\big)
    \end{equation}

\textbf{Step 3]}

Show that

\begin{equation}
	P\Big(  R( f_n)-R(f_a)\geq t\Big) \leq 2|F|\exp\Big( -\frac{nt^2}{2(b-a)^2} \Big)\\
\end{equation}

\textbf{Step 4]}

We write

\begin{equation}
    \delta = 2|F|\exp\Big( -\frac{nt^2}{2(b-a)^2} \Big)
\end{equation}

Show that with probability larger than $ 1-\delta$,

\begin{equation}
    R(f_n)\leq R(f_a)+\sqrt{\frac{2(b-a)^2\Big(\log(\frac{2}{\delta})+\log(|F|)\Big)}{n}}
\end{equation}

In which situations do we have for instance that $a=0$ and $b=1$ ?
\\

{\color{Purple}{In \textbf{tp\_02\_ols}}, we observed a rate in $ \frac{\sigma^2d}{n}$, hence faster than $\mathcal{O}(\frac{1}{\sqrt{n}})$. How can we interpret this ?}

For a more advanced discussion on optimization and statistical bounds, you can have a look at:

\url{https://francisbach.com/scaling-laws-of-optimization/} 
        
\end{document}
